\section{Zadanie 15}

Udowodnij, że jeśli $f, g : \Omega \to \mathbb{C}$ są mierzalne, to
\[
\int_E (f + g) \, d\mu = \int_E f \, d\mu + \int_E g \, d\mu.
\]

\textbf{Wskazówka:}  

Najpierw udowodnij równość dla $f$ i $g$ będących funkcjami schodkowymi.  
Następnie skorzystaj z twierdzenia o aproksymacji funkcji mierzalnej funkcjami schodkowymi.  
Przejście od funkcji nieujemnych do funkcji rzeczywistych i zespolonych jest bardzo proste.

\subsection*{Rozwiązanie}
\textbf{Etap 1: Liniowość całki dla funkcji schodkowych}
\\

Niech $f, g : \Omega \to \mathbb{C}$ będą funkcjami schodkowymi, tzn.
\[
f = \sum_{i=1}^{n} a_i \chi_{A_i}, \qquad
g = \sum_{j=1}^{m} b_j \chi_{B_j},
\]
gdzie $A_i, B_j \subseteq \Omega$ są mierzalnymi i rozłącznymi zbiorami w obrębie $\Omega$, a $a_i, b_j \in \mathbb{C}$.

Wiemy, że
\[
(f+g)(x) = f(x) + g(x).
\]
Ponieważ $f(x) = a_i$ dla $x \in A_i$, a $g(x) = b_j$ dla $x \in B_j$, otrzymujemy:
\[
(f+g)(x) = a_i + b_j \quad \text{dla } x \in A_i \cap B_j.
\]

Zatem funkcja $f+g$ jest również schodkowa i można ją zapisać jako
\[
f + g = \sum_{i=1}^{n} \sum_{j=1}^{m} (a_i + b_j) \chi_{A_i \cap B_j}.
\]

Z definicji całki dla funkcji schodkowej mamy:
\[
\int_E (f+g)\, d\mu = \sum_{i=1}^{n} \sum_{j=1}^{m} (a_i + b_j) \, \mu(E \cap A_i \cap B_j).
\]

Rozdzielając sumę i korzystając z własności sum:
\[
\int_E (f+g)\, d\mu
= \sum_{i=1}^{n} \sum_{j=1}^{m} a_i \, \mu(E \cap A_i \cap B_j)
+ \sum_{i=1}^{n} \sum_{j=1}^{m} b_j \, \mu(E \cap A_i \cap B_j).
\]

Korzystając z rozłączności zbiorów $B_j$ i $A_i$, otrzymujemy:
\[
\sum_{j=1}^{m} \mu(E \cap A_i \cap B_j) = \mu(E \cap A_i), \qquad
\sum_{i=1}^{n} \mu(E \cap A_i \cap B_j) = \mu(E \cap B_j),
\]
czyli
\[
\int_E (f+g)\, d\mu = \sum_{i=1}^{n} a_i \, \mu(E \cap A_i) + \sum_{j=1}^{m} b_j \, \mu(E \cap B_j).
\]

Z definicji całki funkcji schodkowej mamy wprost:
\[
\int_E f \, d\mu = \sum_{i=1}^{n} a_i \, \mu(E \cap A_i), \qquad
\int_E g \, d\mu = \sum_{j=1}^{m} b_j \, \mu(E \cap B_j),
\]
więc wreszcie:
\[
\int_E (f+g)\, d\mu = \int_E f\, d\mu + \int_E g\, d\mu.
\]

\vspace{1em}

\textbf{Etap 2: Liniowość całki dla funkcji mierzalnych nieujemnych}
\\

Z twierdzenia o aproksymacji funkcji mierzalnej funkcjami schodkowymi, weźmy $S_1, \dots, S_n$ i $T_1, \dots, T_n$ takie, żeby
\[
\lim_{n \to \infty} S_n = f, \qquad \lim_{n \to \infty} T_n = g.
\]

Wtedy
\[
\lim_{n \to \infty} (S_n + T_n) = f+g.
\]

Z liniowości całki dla funkcji schodkowych otrzymujemy dla funkcji nieujemnych mierzalnych:
\[
\int_E (f+g) \, d\mu = \lim_{n \to \infty} \int_E (S_n + T_n) \, d\mu 
= \lim_{n \to \infty} \int_E S_n \, d\mu + \lim_{n \to \infty} \int_E T_n \, d\mu 
= \int_E f \, d\mu + \int_E g \, d\mu.
\]

\vspace{1em}

\textbf{Etap 3: Liniowość całki dla funkcji rzeczywistych}
\\

Rozbijmy
\[
f = f_+ - f_-, \qquad g = g_+ - g_-,
\]
gdzie
\[
f_+ = \max\{0,f\}, \quad f_- = -\min\{0,f\},
\]
analogicznie dla funkcji $g$.

Jako że $f_+, f_-, g_+, g_-$ są funkcjami mierzalnymi, nieujemnymi, zastosujemy w dowodzie wynik z Etapu 2:
\[
\begin{aligned}
\int_E (f+g) \, d\mu 
&= \int_E (f_+ - f_-) + (g_+ - g_-) \, d\mu \\
&= \int_E (f_+ + g_+) - (f_- + g_-) \, d\mu \\
&= \int_E (f_+ + g_+) \, d\mu - \int_E (f_- + g_-) \, d\mu \\
&= \int_E f_+ \, d\mu + \int_E g_+ \, d\mu - \int_E f_- \, d\mu - \int_E g_- \, d\mu \\
&= \int_E f \, d\mu + \int_E g \, d\mu.
\end{aligned}
\]

\vspace{1em}

\textbf{Etap 4: Liniowość całki dla funkcji zespolonych}
\\

Niech $f, g : \Omega \to \mathbb{C}$ będą funkcjami mierzalnymi.  
Rozbijamy je na części rzeczywistą i urojoną:
\[
f = \Re(f) + i \Im(f), \qquad g = \Re(g) + i \Im(g),
\]
gdzie $\Re(f), \Im(f), \Re(g), \Im(g)$ są funkcjami rzeczywistymi mierzalnymi.

Zastosowując liniowość całki dla funkcji rzeczywistych (Etap 3), mamy:
\[
\int_E (f+g) \, d\mu = \int_E \big( (\Re(f)+\Re(g)) + i (\Im(f)+\Im(g)) \big) \, d\mu.
\]

Rozdzielamy całkę na część rzeczywistą i urojoną:
\[
\int_E (f+g) \, d\mu = \int_E (\Re(f)+\Re(g)) \, d\mu + i \int_E (\Im(f)+\Im(g)) \, d\mu.
\]

Korzystając z liniowości całki dla funkcji rzeczywistych:
\[
\int_E (\Re(f)+\Re(g)) \, d\mu = \int_E \Re(f) \, d\mu + \int_E \Re(g) \, d\mu, \quad
\int_E (\Im(f)+\Im(g)) \, d\mu = \int_E \Im(f) \, d\mu + \int_E \Im(g) \, d\mu.
\]

Łącząc części rzeczywistą i urojoną, otrzymujemy:
\[
\int_E (f+g) \, d\mu = (\int_E \Re(f) \, d\mu + i \int_E \Im(f) \, d\mu) + (\int_E \Re(g) \, d\mu + i \int_E \Im(g) \, d\mu) = \int_E f \, d\mu + \int_E g \, d\mu.
\]

Zatem liniowość całki zachodzi również dla funkcji zespolonych.
